\uuid{VmzH}
\exo7id{1677}
\titre{exo7 1677}
\auteur{ortiz}
\organisation{exo7}
\datecreate{1999-04-01}
\isIndication{false}
\isCorrection{false}
\chapitre{Réduction d'endomorphisme, polynôme annulateur}
\sousChapitre{Sous-espace stable}
\cpge{mpsi-mp2i.espaces-vectoriels.b, pcsi.espaces-vectoriels.b, ptsi.espaces-vectoriels.a, mp-mpi.reduction-endomorphismes.b, pc.reduction.a, psi.reduction-endomorphismes.a, pt.reduction-endomorphismes.a, mp-mpi.reduction-endomorphismes.d, pc.reduction.c, psi.reduction-endomorphismes.c, pt.reduction-endomorphismes.c, mp-mpi.reduction-endomorphismes.g, pc.complements-espaces-vectoriels.d, psi.complements-espaces-vectoriels.d, bcpst2.valeurs-propres-vecteurs-propres.a, bcpst2.valeurs-propres-vecteurs-propres.b}
\module{Algèbre}
\niveau{L2}
\difficulte{}

\contenu{
\texte{
Soient $f\in\mathcal{L}(E)$ et $a\in E.$
}
\begin{enumerate}
    \item \question{Montrer que le plus petit sous-espace vectoriel de $E$ contenant $a$ et
stable par $f $ est
$F_a=\text{vect}\left\{f^k(a):k\in\Nn\right\}$.}
    \item \question{Montrer que si $\dim(E)=n$ alors $F_a=\text{vect}\left\{f^k(a):k=0,...,n-1\right\}.$}
    \item \question{Soit l'endomorphisme $f$ de $\Rr^3$ canoniquement associ\'e \`a la
matrice $A=\left( \begin{smallmatrix}
1&1&-1\\
-2&-2&2\\
1&1&-1
\end{smallmatrix}\right) \in \mathcal{M}_3(\Rr)$.
Montrer qu'il n'existe pas $a\in \Rr^3$ tel que
$F_a=\Rr^3.$ G\'en\'eraliser \`a un endomorphisme
diagonalisable.}
\end{enumerate}
}
